% !TeX root = main.tex
% Это файл-вставка без преамбулы: собирается только через main.tex.
\section{Анализ задачи визуализации поиска по дереву методом Монте-Карло и методов семантического масштабирования}

Настоящая работа продолжает выпускную квалификационную работу бакалавра [1], в рамках которой в программной среде GGP Base Package был разработан и интегрирован модуль визуализации построения дерева поиска методом Монте-Карло. Созданный модуль обеспечил наблюдение за ростом дерева, однако его применимость ограничена сравнительно небольшими деревьями. В настоящем разделе приводятся теоретические сведения о методе поиска и порождаемых им данных, рассматриваются особенности дерева поиска как объекта визуализации и задачи его анализа, описывается существующее программное решение, анализируются методы масштабирования больших иерархий и реализующие их программные средства, после чего формулируются проблема, цель работы, исследовательские вопросы и требования к разрабатываемому решению.

\subsection{Поиск по дереву методом Монте-Карло}

\subsubsection{Основные понятия теории игр и интеллектуальных агентов}

Игра — это математическая модель стратегической ситуации принятия решения (конфликт, сотрудничество и тому подобное), включающая взаимодействие. Во взаимодействие вовлечены различные действующие лица, так называемые игроки, которые совместно определяют результат, и каждый старается получить результат, который является для него наилучшим [2].

Игра описывается следующим набором компонент [3, 4]:

\begin{equation}
G = \langle R,\; S,\; s_0,\; A,\; noop,\; terminal(\cdot),\; legal_i(\cdot),\; next(\cdot),\; payoff(\cdot) \rangle
\end{equation}

где R = \{1, …, r\} — набор из r ролей. Роль — это субъект, который совершает действия в игре.

S — множество всех возможных состояний игры. Под состоянием игры понимается множество устойчивых значений переменных параметров игры.

s₀ ∈ S — начальное состояние игры, то есть состояние игры в самом её начале до совершения действий игроками.

terminal: S → \{true, false\} — функция проверки каждого состояния на терминальность. Терминальное состояние, или исход, — состояние игры, которое принимается в момент завершения игры и из которого каждый игрок не имеет легальных действий.

A = \{A₁, …, A\_r\} — множество действий для каждой роли, которое может быть совершено по крайней мере в одном из состояний игры. Действие — это то, что совершает конкретная роль, меняя состояние игры.

noop — специальное действие, обозначающее пустое действие.

legalᵢ: S → P(Aᵢ) — функция, определяющая легальность действий в каждом состоянии для i-й роли, где P(Aᵢ) — множество всех подмножеств множества Aᵢ. Легальное действие — это действие, которое может быть совершено данным игроком в некотором состоянии игры.

next: S × legal(s, 1) × legal(s, 2) × … × legal(s, r) → S — функция, определяющая переход в следующее состояние игры.

payoff: R × S → ℝʳ — функция получения набора платежей для всех ролей в терминальном состоянии игры. Платёж — величина, которая характеризует предпочтения данного игрока в заданной игре.

Таким образом, перечисленные компоненты порождают дерево игры, узлами которого являются состояния из множества S, а рёбрами — действия игроков, приводящие к переходам между состояниями согласно функции next(·). Корневой узел соответствует начальному состоянию s₀, а листовые узлы — терминальным состояниям, для которых функция terminal(·) возвращает true.

Игрок — это рациональный индивид, который принимает решения в игре, представляя некоторую роль в игре, и имеет заинтересованность в наибольшем выигрыше. Рациональность игрока означает то, что каждый игрок совершает такие действия, которые приводят к наибольшему выигрышу в каждый момент игры, в котором данный игрок действует [41].

Цель игрока — достичь исхода с максимально возможным платежом при условии, что остальные игроки также стремятся к своим предпочтительным исходам [42]. При этом предпочтительные исходы всех игроков не обязательно совпадают. Для достижения цели игрок должен реализовать в игре, точнее в матче, свою оптимальную стратегию.

Стратегия — это полный план действий, который задаёт для каждого состояния игры действие, которое этот игрок намеревается выбрать, если окажется в этом состоянии, то есть strategyᵢ: S → Aᵢ — стратегия i-го игрока [2].

Равновесие Нэша — концепция, согласно которой ни один из игроков не может улучшить свой результат, изменив свою стратегию при условии, что стратегии остальных игроков остаются неизменными [43].

В работе рассматриваются детерминированные игры с полной информацией — игры, в которых все игроки имеют доступ ко всей информации о состоянии игры, например шахматы [44].

Профиль стратегий — множество стратегий всех игроков. Профиль однозначно определяет исход матча и представляет в дереве игры путь от корня до листа. Для детерминированных игр с полной информацией профиль, содержащий оптимальные стратегии всех игроков, удовлетворяет концепции равновесия Нэша в подиграх [45].

Интеллектуальный агент — некоторая программная сущность, способная действовать в интересах достижения целей, поставленных перед ним владельцем или пользователем. Обычно предполагается, что интеллектуальный агент [3, 4]:

\begin{itemize}

  \item действует автономно или совместно с другими компьютерными и интеллектуальными системами;

  \item выполняет шаблонные предписанные действия и действия, требующие активности и учёта состояния окружающей среды;

  \item может генерировать цели и действовать рациональным образом для их достижения;

  \item в той или иной степени способен к обучению и корпоративным действиям, мобилен.

\end{itemize}

\subsubsection{Алгоритм поиска по дереву методом Монте-Карло}

Поиск по дереву методом Монте-Карло (Monte Carlo Tree Search) — эвристический алгоритм принятия решений в задачах, допускающих представление деревом состояний [5]. Метод итеративно наращивает лишь фрагмент дерева игры, оценивая перспективность узлов по статистике случайных симуляций, доигрывающих партию до завершения: развиваются преимущественно те направления, которые алгоритм считает перспективными.

Каждая итерация складывается из четырёх этапов [6]. На этапе выбора (selection) алгоритм, начиная с корня, последовательно переходит к дочернему узлу с наибольшим значением критерия отбора, пока не достигнет листа. На этапе расширения (expansion) лист дополняется одной или несколькими дочерними узлами, отвечающими допустимым продолжениям. На этапе симуляции (simulation) из добавленного узла партия случайно доигрывается до терминального состояния. На этапе обратного распространения (backpropagation) полученный результат переносится по пройденному пути от листа к корню, корректируя в каждом узле накопленную ценность и счётчик посещений. Последовательность фаз одной итерации показана на рисунке 1.1.

\gostfigure{../../images/chapter1/fig-1-1-phases.png}{14cm}{Этапы итерации поиска по дереву методом Монте-Карло}

Чтобы уравновесить исследование ещё не изученных участков дерева и использование уже оценённых, в классическом варианте метода применяют критерий UCT (Upper Confidence Bound for Trees): на этапе выбора предпочтение получает дочерний узел j, максимизирующий величину

\begin{equation}
UCT_j = \bar{X}_j + C\sqrt{\dfrac{\ln N}{n_j}}
\end{equation}

где X̄\_j — средняя награда узла j, N — число посещений её родителя, n\_j — число посещений самого узла, C регулирует соотношение между исследованием и использованием [7]. Первое слагаемое поощряет узлы с высокой средней оценкой, второе возрастает для редко посещаемых.

Характерной особенностью метода является асимметричный рост дерева: перспективные ветви развиваются глубже и шире, малоперспективные остаются слаборазвитыми. Число итераций на практике ограничивается заданным значением либо отведённым временем. Накапливаемая в узлах статистика — посещения, средние оценки и число использований каждого действия — и является предметом визуализации, поскольку отражает, какие участки пространства поиска алгоритм счёл перспективными и как распределил между ними ресурсы.

\subsection{Дерево поиска как объект визуализации}

\subsubsection{Особенности дерева поиска}

Из природы алгоритма вытекает ряд особенностей дерева поиска, определяющих требования к его визуализации. Прежде всего это размер. В наблюдаемых партиях игры «Четыре в ряд» дерево одного хода при 1200 итерациях насчитывает около 2500 узлов, журнал полной партии фиксирует порядка 42 тысяч итераций, а при десятках секунд на ход счёт итераций доходит до миллионов. Уже при секундах на ход число узлов превышает число элементов, одновременно различимых на экране по горизонтали, поэтому часть сведений обязана оставаться скрытой.

Распределение посещений по узлам при этом резко неоднородно: вследствие асимметричного роста основную массу посещений сосредотачивает небольшое число доминирующих ветвей, тогда как множество остальных остаётся слаборазвитым. Крайняя форма этой неоднородности наблюдалась на почти решённой позиции: ход, на который алгоритм затратил 6,2 млн итераций, дал дерево всего из 713 узлов — практически весь объём поиска пришёлся на одну узкую глубокую ветвь, то есть число узлов не растёт пропорционально числу итераций. Визуализация обязана выделять доминирующие ветви, не теряя из виду слабые: именно слабые ветви показывают, какие варианты алгоритм рассмотрел и отверг.

Дерево не является статичным: оно растёт от итерации к итерации, а при переходе к следующему ходу подвергается обрезке и смене корня. Анализ требует наблюдения за процессом изменения дерева, а не только за итоговым состоянием, поскольку многие особенности работы алгоритма проявляются на ранних итерациях и маскируются после их большого числа [1].

В рассматриваемой реализации дерево переиспользуется между ходами игры: счётчики посещений узлов-предков накапливаются на протяжении всей партии, тогда как поддеревья перестраиваются. Вследствие этого абсолютные значения счётчиков посещений несопоставимы между уровнями дерева, и любые доли необходимо вычислять относительно соседних узлов одного родителя, а не относительно корня. Это свойство данных существенно влияет на корректность кодирования размеров узлов при визуализации.

Реализация метода выполнена в программной среде GGP Base Package — платформе общего игрового искусственного интеллекта, позволяющей агенту играть в произвольные игры, описанные на языке GDL, без предварительного знания их правил [8]. Поскольку набор игр не фиксирован и не известен заранее, визуализация не может опираться на игровые представления — отрисовку доски, фигур или иных специфических для игры элементов. Допустимым предметом отображения остаётся только структура дерева поиска и статистика его узлов, то есть данные, общие для всех игр. Требование независимости от конкретной игры исключает из рассмотрения класс решений, привязанных к отдельным играм, и направляет работу в сторону универсальных методов визуализации деревьев.

\subsubsection{Задачи анализа работы алгоритма}

Требования к визуализации определяются тем, какие вопросы исследователь задаёт дереву поиска. Эти вопросы возникают при отладке реализации, при настройке параметров и при сравнении вариаций метода и сводятся к семи типовым:

\begin{enumerate}

  \item какой ход выбран и насколько уверенно — нужно видеть распределение посещений между ходами в корне и отрыв лидера;

  \item когда предпочтение сменилось — ранние итерации нередко выводят в лидеры ход, впоследствии отвергаемый, и момент смены характеризует устойчивость решения;

  \item какие продолжения рассмотрены и отвергнуты — отвергнутые ветви показывают, что именно алгоритм счёл бесперспективным, и их исчезновение с экрана равносильно потере половины сведений о поиске;

  \item что произошло на отдельной итерации — какой лист выбран, какие узлы созданы, откуда выполнялась симуляция и по какому пути распространён результат;

  \item где поиск сосредоточен сейчас, а не за всё время хода;

  \item чем различаются два прогона, выполненные разными вариациями метода или с разным временем на ход;

  \item встречалась ли одна позиция в разных ветвях дерева.

\end{enumerate}

Перечисленные вопросы объединяет то, что ни один из них не решается ни общим обзором, ни детальным просмотром по отдельности: каждый требует перехода между масштабами — от распределения усилий по всему дереву к статистике отдельного узла и обратно. Именно поэтому задача сводится не к отображению дерева, а к навигации по нему.

Потребность в таком анализе не ограничивается отладкой: дерево поиска служит и источником обоснования выбранного хода. В работах по объяснимости поиска по дереву методом Монте-Карло объяснение строится как интерактивное исследование дерева с обобщением поддеревьев и сопоставлением выбранного действия с отвергнутыми [22], что предъявляет к визуализации те же требования.

\subsubsection{Противоречие объёма данных и экранного пространства}

Основное противоречие визуализации больших деревьев состоит в несоответствии между объёмом данных и ограниченностью экранного пространства. Число узлов дерева поиска превышает число доступных по горизонтали пикселей, поэтому попытка отобразить дерево целиком приводит либо к информационной перегрузке, при которой узлы накладываются друг на друга и подписи становятся нечитаемыми, либо к потере детализации, при которой узлы вырождаются в неразличимые точки. Простое геометрическое уменьшение масштаба не разрешает это противоречие: оно пропорционально уменьшает все элементы, не учитывая их значимости, в результате чего важные и второстепенные узлы становятся одинаково нечитаемыми.

Общепринятым способом организации навигации по большим наборам данных является принцип последовательного перехода от обзора к деталям: сначала предъявляется общий обзор, затем средства масштабирования и фильтрации, и лишь по явному запросу — детальная информация об отдельном элементе [9]. Применительно к дереву поиска это означает, что на обзорном масштабе видны общая форма дерева и его доминирующие ветви, при приближении — детали интересующего поддерева, а полная статистика узла доступна по запросу. Реализацию осложняет то, что дерево не помещается на экран целиком, распределение посещений неоднородно, а состав отображаемых узлов должен изменяться при приближении, не нарушая устойчивости картины. Перечисленное и определяет содержание задачи.

\subsection{Существующее программное решение}

Модуль, разработанный в предшествующей работе [1], представляет дерево в виде узлов-кругов, размещаемых слева направо, с размером, пропорциональным количеству посещений, и цветовой кодировкой относительной интенсивности исследования; динамика отображается пошаговым воспроизведением снимков дерева. Этот подход успешно решает задачу наблюдения за ростом дерева на малых масштабах, но обладает ограничениями, проявляющимися при увеличении размера дерева. Представление узлов-кругов не сохраняет читаемость при числе узлов порядка нескольких тысяч: круги либо перекрываются, либо вырождаются в точки, а соединяющие их рёбра образуют визуальный шум. Управление детализацией возложено на пользователя — для упрощения картины требуется вручную задавать ограничения по глубине и минимальному проценту посещений, что перекладывает на исследователя задачу, которую система должна решать автоматически, и требует заранее знать, какую часть дерева следует скрыть. Кроме того, при сворачивании и разворачивании ветвей геометрия перестраивается и узлы смещаются, вследствие чего теряется пространственный контекст: модуль решает задачу отображения, но не решает задачи навигации по большому дереву. Его внешний вид приведён на рисунке 1.2.

\gostfigure{../../images/chapter1/fig-1-2-existing.png}{14cm}{Визуализация дерева поиска в существующем модуле}

Порядок работы исследователя с существующим модулем показан на рисунке 1.3. Чтобы получить читаемую картину, пользователь вынужден подбирать ограничения по глубине и минимальной доле посещений вручную и повторять подбор после каждого изменения дерева — цикл «настроить фильтры → оценить картину → перенастроить» замыкается на исследователя, а не на систему.

\gostfigure{../../images/chapter1/fig-1-3-activity.png}{9.8cm}{Процесс анализа дерева поиска с использованием существующего модуля (диаграмма деятельности)}

Таким образом, уровень детализации здесь задаётся вручную, а его изменение нарушает устойчивость картины, что и обосновывает обращение к методам семантического масштабирования.

\subsection{Методы масштабирования больших иерархий}

Сформулированное противоречие указывает направление решения, но не определяет конкретного представления дерева. В настоящем подразделе излагаются принципы, связывающие масштаб отображения с уровнем детализации, рассматриваются применяемые для больших иерархий представления — узло-связевые и пространственно-заполняющие, — а также способы обращения с узлами, которые на экране не помещаются, после чего выполняется выбор базового представления.

\subsubsection{Семантическое и геометрическое масштабирование}

Семантическое масштабирование — это метод взаимодействия, при котором изменение масштаба отображения меняет не только размер элементов, но и состав и форму представления отображаемой информации [10, 11]. В отличие от геометрического масштабирования, при котором объекты лишь пропорционально увеличиваются или уменьшаются, при семантическом масштабировании один и тот же объект на разных уровнях приближения может отображаться качественно по-разному: на обзорном масштабе — обобщённо или скрыто, на детальном — с полной информацией. Понятие введено в работах по интерфейсам с непрерывным масштабированием и получило развитие в обзоре интерфейсов совмещения обзора и детали, масштабирования и совмещения фокуса и контекста [10, 11, 12].

Применительно к дереву поиска семантическое масштабирование означает, что уровень детализации отображаемого поддерева определяется текущим масштабом: на обзорном масштабе показываются только крупные, доминирующие ветви, а мелкие узлы скрываются; при приближении к выбранной области ранее скрытые узлы постепенно проявляются. Это позволяет удерживать плотность отображаемой информации приблизительно постоянной независимо от размера дерева.

\subsubsection{Совмещение фокуса и контекста и степень интереса}

С семантическим масштабированием тесно связаны два принципа информационной визуализации. Принцип совмещения фокуса и контекста предполагает одновременное отображение детально рассматриваемой области и её окружения в общем поле зрения, что позволяет пользователю не терять ориентацию при углублении в детали [12]. Принцип степени интереса, предложенный Фёрнасом в модели обобщённых представлений типа «рыбий глаз» [13], формализует идею о том, что для каждого элемента можно вычислить численную меру его значимости для пользователя, зависящую от внутренней важности элемента и его удалённости от текущего фокуса; отображаются лишь элементы с достаточно высокой степенью интереса.

Для дерева поиска естественной мерой внутренней важности узла служит его доля посещений, а удалённость от фокуса определяется текущей областью приближения. Сочетание этих принципов даёт основу для автоматического, а не задаваемого вручную, управления детализацией, что устраняет основной недостаток существующего модуля.

\subsubsection{Уровни детализации и картографическая генерализация}

Наиболее разработанной областью применения семантического масштабирования являются картографические интерфейсы: при отдалении карта показывает только крупные объекты, при приближении проявляются всё более мелкие, причём объекты и подписи не уменьшаются пропорционально, а появляются и исчезают. В картографии переход между уровнями детализации описан процедурами генерализации [21]: отбором объектов, сохраняемых на данном масштабе, агрегированием мелких объектов в обобщённые, упрощением форм и снятием подписей.

Каждой из этих процедур соответствует операция над деревом поиска. Отбору отвечает скрытие узлов, не достигающих порога значимости на текущем масштабе; агрегированию — представление скрытых поддеревьев обобщёнными участками, сохраняющими суммарную статистику; снятию подписей — сокращение состава подписи узла по мере уменьшения его экранного размера. Такое соответствие показывает, что задача визуализации дерева поиска структурно повторяет задачу картографической генерализации, и связь уровня детализации с масштабом, отработанная в картографии, переносится на дерево поиска. Из рассмотренных принципов только семантическое масштабирование обеспечивает одновременно автоматическое управление детализацией, устойчивость положения элементов и сохранение контекста: методики, основанные на искажении пространства, жертвуют устойчивостью, раскрытие по степени интереса ставит положение элементов в зависимость от выбранного фокуса, а ручная фильтрация не даёт ни одного из трёх свойств. Поэтому семантическое масштабирование и положено в основу предлагаемого подхода.

\subsubsection{Узло-связевые представления}

В узло-связевых представлениях узлы изображаются геометрическими фигурами, а рёбра — соединяющими их линиями; такое представление использовано и в существующем модуле [1]. Для повышения масштабируемости подобных представлений разработан ряд методов.

Гиперболический браузер размещает дерево на гиперболической плоскости, отображаемой в круг: узел в фокусе занимает центр и изображается крупно, а удалённые узлы экспоненциально сжимаются к границе круга [14]. Благодаря экспоненциальному сжатию метод вмещает деревья из тысяч узлов и служит классическим примером совмещения фокуса и контекста. Однако размер узла в нём определяется расстоянием до фокуса, а не собственной метрикой узла, поэтому неоднородность посещений этим представлением не передаётся; масштабирование сводится к дискретной смене фокуса, а не к непрерывному изменению уровня детализации; при каждой смене фокуса все узлы приходят в движение, а искажение геометрии затрудняет и удержание пространственного контекста, и чтение путей «предок — потомок».

Метод SpaceTree сочетает обычную компоновку дерева с динамическим управлением раскрытием ветвей: поддеревья, не помещающиеся на экран, сворачиваются в пиктограммы, размер которых кодирует количество скрытых узлов [15]. Деревья со степенью интереса применяют модель Фёрнаса для автоматического выбора отображаемых узлов, агрегируя малозначимые ветви [16]. Эти методы реализуют переход от обзора к деталям и сохраняют привычную читаемость путей, однако плохо передают резко неоднородное распределение количественной метрики, поскольку размер узла в них слабо связан с его значимостью, а при сворачивании и разворачивании ветвей компоновка перестраивается и узлы смещаются — пространственная устойчивость обеспечивается лишь частично.

\subsubsection{Пространственно-заполняющие представления}

В пространственно-заполняющих представлениях иерархия кодируется вложенностью или смежностью областей, а не явными рёбрами, что обеспечивает плотное использование экранного пространства. Ярусная диаграмма (icicle plot) изображает дерево рядами прямоугольников: каждый уровень глубины образует строку, ширина прямоугольника пропорциональна некоторой метрике узла, а дочерние узлы располагаются под родительским, деля его ширину [17]. Отношение «родитель — потомок» считывается из взаимного расположения строк без явных рёбер. Флейм-граф (flame graph) представляет собой ярусную диаграмму, в которой ширина прямоугольника кодирует количественную метрику; метод широко применяется для анализа профилей производительности и наглядно передаёт доминирующие по весу пути [18]. Диаграмма вложенных прямоугольников (treemap) кодирует иерархию площадью вложенных прямоугольников, пропорциональной метрике узла [19]; она удобна для отображения распределения величины по листьям, но хуже передаёт глубину и структуру путей, а приближение к выбранной области требует перекладки прямоугольников, при которой их взаимное расположение меняется, — с одноосевым масштабированием площадная компоновка сочетается плохо. Радиальным аналогом ярусной диаграммы является диаграмма sunburst, в которой уровни глубины образуют концентрические кольца [20]; с ростом глубины кольца утончаются, чтение долей по дугам менее точно, чем по отрезкам, а растяжение одной оси в радиальной геометрии невозможно без смены системы координат. Примеры названных представлений приведены на рисунке 1.4.

\gostfigure{../../images/chapter1/fig-1-4-layouts.png}{14cm}{Пространственно-заполняющие представления дерева: ярусная диаграмма (флейм-граф), диаграмма вложенных прямоугольников и диаграмма sunburst}

Диаграмма вложенных прямоугольников применялась и непосредственно к деревьям поиска. В работе [23] этим способом визуализированы деревья альфа-бета-поиска шахматных программ: площадь прямоугольника пропорциональна числу узлов поддерева, а цвет кодирует глубину узла, отсечения и ветви, просмотренные напрасно из-за неудачного порядка ходов. Эта работа является ближайшим предшественником настоящей по предмету визуализации, и выявленные в ней затруднения показательны. Поочерёдное деление области то по вертикали, то по горизонтали даёт крайне вытянутые прямоугольники, поскольку за отсечением обычно следует полное развёртывание очередного уровня; пропорции исправляет выравнивающая раскладка [24], но она упорядочивает прямоугольники по величине и теряет порядок ходов, существенный для чтения формы дерева. Предложенное в [23] решение — выравнивающая раскладка без упорядочивания по величине, сохраняющая порядок ценой менее правильных пропорций. Второе затруднение существеннее: при приближении или изменении размера окна меняются пропорции области, а вместе с ними перестраивается вся раскладка; единственным средством против этого в [23] служит ручная блокировка раскладки, вновь ухудшающая пропорции. Оба затруднения порождены делением площади по двум измерениям, в котором порядок узлов и пропорции прямоугольников соперничают за одни и те же степени свободы. Кроме того, дерево в [23] неизменно во времени — оно читается из файла записанной партии, — а само представление привязано к шахматам.

Преимущество ярусной диаграммы и флейм-графа для рассматриваемой задачи состоит в том, что ширина прямоугольника напрямую кодирует количественную метрику узла, благодаря чему доминирующие по посещениям ветви визуально доминируют и по занимаемой ширине. Это свойство точно соответствует асимметричной природе дерева поиска. Кроме того, прямоугольная компоновка естественно сочетается с горизонтальным масштабированием: приближение реализуется как растяжение горизонтальной оси, при котором узкие прямоугольники расширяются и становятся доступными для детального рассмотрения, — по этой схеме устроены интерактивные просмотрщики флейм-графов в средствах анализа производительности [18]. Вместе с тем ни одно из существующих средств не предназначено для деревьев поиска и не учитывает их особенностей — переиспользования дерева между ходами, необходимости различать объём поиска и его результат, наблюдения за ростом во времени и независимости от конкретной игры.

\subsubsection{Обращение с узлами, не помещающимися на экран}

Перечисленные представления отвечают на вопрос, как разместить узлы на экране, но не на вопрос, что делать с узлами, которые на нём не помещаются. Этому посвящено отдельное направление — иерархическая агрегация [25], в которой группы элементов заменяются обобщёнными элементами. Направление формулирует требования к такой замене: обобщённый элемент должен нести краткую сводку о том, что он объединяет, визуально отличаться от обычного элемента, допускать развёртывание и не искажать существенных свойств данных, а число одновременно отображаемых элементов ограничивается заданным бюджетом, который удерживает одновременно и загромождённость изображения, и стоимость его построения. Применительно к дереву поиска отсюда следует, что скрытые поддеревья обязаны оставаться на экране обобщёнными участками с собственной статистикой, а не просто исчезать.

Близкое по смыслу направление — управление уровнем детализации, при котором состав отображаемых элементов определяется их размером на экране, а пространственно близкие узлы и рёбра объединяются; это позволяет интерактивно отображать графы с миллионами элементов [26]. Такие методы решают вычислительную сторону масштабируемости и обосновывают сам приём отбора по экранному размеру, однако не придают объединённым элементам содержательного смысла — для дерева поиска его приходится определять отдельно.

Наконец, постепенная выдача [27] рассматривает случай, когда данные или представление выдаются частями: сначала грубое приближение, затем уточнение. Обзор этого направления выделяет иерархические уровни обобщения как естественный механизм уточнения, что согласуется с построением, при котором приближение к ветви догружает её подробности.

\subsubsection{Сравнение методов и выбор базового представления}

Для выбора базового представления выделены пять критериев: масштабируемость, то есть сохранение читаемости при тысячах узлов; способность кодировать неоднородную метрику; совместимость с семантическим масштабированием; пространственная устойчивость, то есть отсутствие смещения узлов при изменении детализации; читаемость структуры путей. Пространственная устойчивость включена в число критериев не как самоочевидное благо: экспериментальные сравнения анимации и серий неподвижных изображений показывают, что выигрыш от сохранения взаимного расположения элементов зависит от решаемой задачи и наиболее заметен там, где требуется повторно находить ранее рассмотренный элемент и прослеживать пути [28], — а именно эти действия и составляют работу с деревом поиска. Результаты сравнения методов по этим критериям сведены в таблицу 1.1.

\begin{gosttable}{}{Сравнение методов визуализации больших деревьев по критериям применимости к визуализации дерева поиска}

\begin{tabularx}{\linewidth}{>{\raggedright\arraybackslash}Xccccc}
\hline
Метод & Масштабируемость & Кодирование неоднородной метрики & Совместимость с масштабированием & Пространственная устойчивость & Читаемость структуры путей \\
\hline
Гиперболический браузер & + & − & ± & − & ± \\
\hline
SpaceTree, деревья со степенью интереса & + & − & + & ± & + \\
\hline
Диаграмма вложенных прямоугольников & + & + & ± & ± & − \\
\hline
Диаграмма sunburst & ± & + & ± & ± & ± \\
\hline
Ярусная диаграмма (флейм-граф) & + & + & + & + & + \\
\hline
\end{tabularx}

\end{gosttable}

Условные обозначения: «+» — критерий выполняется, «±» — выполняется частично, «−» — не выполняется.

Сравнение показывает, что наилучшим сочетанием свойств обладает ярусная диаграмма с горизонтальным семантическим масштабированием: она масштабируема, кодирует неоднородную метрику посещений шириной прямоугольников и при надлежащем построении координат обеспечивает пространственную устойчивость. Существенно, что деление единственной оси, в отличие от деления площади по двум измерениям [23], не ставит порядок узлов и пропорции прямоугольников в зависимость друг от друга: порядок ходов сохраняется по построению, а изменение масштаба не перестраивает раскладку, поэтому устойчивость не требует ни ручной блокировки, ни ухудшения пропорций. Это представление выбрано в качестве основы для разрабатываемой модели.

\subsection{Обзор существующих программных средств}

Рассмотренные методы определяют, что в принципе возможно при визуализации больших иерархий. Настоящий подраздел устанавливает, насколько эти возможности реализованы в доступных средствах и в какой мере они покрывают задачи анализа, перечисленные в 1.2.2. В обзор включены средства с публично доступной реализацией или сетевой демонстрацией, поддерживающие отображение деревьев поиска либо деревьев родственных алгоритмов и снабжённые документацией или примерами применения; средства, привязанные к единственной игре, не исключались, поскольку именно они составляют основную часть доступных решений. Сравнение выполнено по методике и с результатами, опубликованными в работе [29], и дополнено сопоставлением по свойствам, существенным для навигации по большому дереву.

\subsubsection{Критерии сравнения}

Для сравнения сформулированы двенадцать критериев, отражающих потребности отладки и исследования метода поиска. Критериям присвоены веса, в сумме составляющие сто процентов; наибольший вес получили универсальность по отношению к играм, полнота отображаемых метрик, поддержка промежуточных состояний дерева и интерактивность, поскольку именно они определяют пригодность средства для анализа поведения алгоритма. Критерии и их веса приведены в таблице 1.2.

\begin{gosttable}{}{Критерии сравнения средств визуализации дерева поиска}

\begin{tabularx}{\linewidth}{>{\raggedright\arraybackslash}Xccc}
\hline
Обозначение & Критерий & Вес, \% & Содержание \\
\hline
k1 & Полнота метрик & 15 & Отображение числа посещений, средней оценки и производных от них величин \\
\hline
k2 & Промежуточные состояния дерева & 15 & Отображение дерева на разных этапах его построения, а не только итогового \\
\hline
k3 & Универсальность & 18 & Поддержка произвольных игр и платформ \\
\hline
k4 & Масштабируемость & 8 & Работоспособность при числе узлов свыше тысячи \\
\hline
k5 & Интерактивность & 14 & Масштабирование, панорамирование, фильтрация, контекстные меню \\
\hline
k6 & Воспроизведение построения & 12 & Пуск, пауза и шаг с управляемой скоростью \\
\hline
k7 & Плавность переходов & 5 & Непрерывность визуальных переходов между состояниями дерева \\
\hline
k8 & Качество интерфейса & 5 & Общее впечатление и удобство работы \\
\hline
k9 & Лёгкость освоения & 3 & Трудоёмкость начала работы, интуитивность \\
\hline
k10 & Кодирование величин & 2 & Использование цвета, размера и формы для передачи метрик узла \\
\hline
k11 & Документация & 2 & Наличие руководств, примеров и справки \\
\hline
k12 & Поддержка разных пользователей & 1 & Пригодность для начинающего, студента, исследователя \\
\hline
\end{tabularx}

\end{gosttable}

Интегральная оценка средства вычисляется как взвешенная сумма оценок по всем критериям, нормированная на максимально достижимое значение:

\begin{equation}
E(S) = \dfrac{1}{s_{\max} \cdot 100}\sum_{i=1}^{12} score(S,\, k_i) \cdot w_i
\end{equation}

где S — оцениваемое средство; score(S, k\_i) — экспертная оценка средства S по критерию k\_i по шкале от нуля до пяти; w\_i — вес критерия k\_i в процентах; s\_max = 5 — наибольшая оценка шкалы. Нормировка приводит интегральную оценку к отрезку от нуля до единицы. Оценивание выполнено двумя экспертами, имеющими опыт разработки агентов на основе поиска по дереву методом Монте-Карло в рассматриваемой программной среде; итоговые значения получены усреднением их оценок, формальные меры согласованности не вычислялись.

\subsubsection{Рассмотренные средства}

Рассмотрены десять средств, которые по характеру решаемой задачи образуют четыре группы.

Первую группу составляют учебные визуализаторы: mcts-viz [30], пошаговый разбор AlphaZero и поиска по дереву [31], AlphaZero Tic-Tac-Toe with Monte Carlo Tree Search [32], AlphaZero-Edu [35] и Playing with Monte-Carlo Tree Search [37]. Они наглядно показывают устройство алгоритма: разделяют этапы итерации, подписывают значения критерия отбора у каждого дочернего узла, позволяют выполнять итерации по одной. Все они построены вокруг единственной игры, а размер отображаемого дерева ограничен несколькими десятками узлов: представление рассчитано на разбор отдельных шагов, а не на анализ длительного поиска.

Вторую группу образуют универсальные библиотеки визуализации — D3 [34] и D3-Graphviz [36]. Не завися от предметной области, они поддерживают произвольные деревья и развитое взаимодействие, но не содержат ни понятий метода поиска, ни наблюдения за изменением структуры: связующий код, отбор отображаемых узлов и устойчивость раскладки остаются на исследователе.

Третью группу составляют средства просмотра итогового дерева — Gloqo [33] и просмотр дерева, выгруженного средой GGP Base Package, в виде развёртываемых таблиц [38, 39]. Такой просмотр точно передаёт все хранимые величины и выдерживает большой объём, но зрительного представления дерева не образует: форма дерева и распределение посещений между ветвями остаются невидимыми, а динамика построения недоступна в принципе.

Четвёртую группу составляют программные библиотеки, реализующие сам метод поиска и снабжённые вспомогательными средствами отображения, — MCTS.jl [40]. Они независимы от конкретной игры, однако отображение в них является побочной возможностью: оно ограничено выводом структуры дерева и не предназначено для интерактивного анализа.

\subsubsection{Результаты сравнения}

Оценки средств по критериям и полученные по формуле (1.3) интегральные оценки приведены в таблице 1.3.

\begin{gosttable}{}{Оценки средств визуализации дерева поиска}

\begin{tabularx}{\linewidth}{>{\raggedright\arraybackslash}Xccccccccccccc}
\hline
Средство & k1 & k2 & k3 & k4 & k5 & k6 & k7 & k8 & k9 & k10 & k11 & k12 & E(S) \\
\hline
mcts-viz & 5 & 4 & 0 & 2 & 2 & 4 & 2 & 2 & 4 & 2 & 0 & 2 & 0,530 \\
\hline
Josh Varty & 4 & 4 & 1 & 2 & 2 & 3 & 4 & 3 & 4 & 3 & 3 & 3 & 0,560 \\
\hline
AlphaZero TTT & 4 & 4 & 2 & 3 & 4 & 3 & 3 & 4 & 4 & 3 & 4 & 1 & 0,668 \\
\hline
Gloqo & 4 & 3 & 2 & 3 & 3 & 2 & 3 & 3 & 3 & 3 & 3 & 2 & 0,568 \\
\hline
D3 & 3 & 2 & 5 & 4 & 4 & 1 & 3 & 4 & 3 & 4 & 5 & 1 & 0,656 \\
\hline
AlphaZero-Edu & 3 & 2 & 4 & 3 & 3 & 1 & 2 & 3 & 3 & 2 & 3 & 2 & 0,542 \\
\hline
D3-Graphviz & 3 & 2 & 4 & 4 & 3 & 1 & 2 & 4 & 2 & 4 & 5 & 1 & 0,576 \\
\hline
Playing with MCTS & 3 & 3 & 1 & 2 & 3 & 2 & 2 & 3 & 3 & 2 & 2 & 1 & 0,466 \\
\hline
GGP + JSON Grid & 5 & 1 & 3 & 5 & 0 & 1 & 0 & 1 & 2 & 3 & 0 & 3 & 0,432 \\
\hline
MCTS.jl & 2 & 1 & 5 & 3 & 2 & 1 & 2 & 3 & 2 & 2 & 3 & 1 & 0,482 \\
\hline
\end{tabularx}

\end{gosttable}

Распределение оценок показывает устойчивую закономерность: высокие значения по полноте метрик и промежуточным состояниям достигаются лишь средствами, привязанными к единственной игре, а высокие значения по универсальности — средствами, не содержащими понятий метода поиска. Ни одно средство не получает высоких оценок одновременно по критериям k1–k6. Наибольшая интегральная оценка составляет 0,668, наименьшая — 0,432. Модуль, разработанный в предшествующей работе, получает по той же методике оценку 0,752 [29] за счёт сочетания независимости от игры, воспроизведения построения и фильтрации, но ограничен деревьями в несколько тысяч узлов и ручным управлением детализацией (1.3).

\subsubsection{Сопоставление по свойствам, существенным для навигации}

Приведённые критерии описывают средство как инструмент отладки в целом и не выделяют свойств, от которых зависит работа с большим деревом. Поэтому группы средств дополнительно сопоставлены по пяти свойствам, вытекающим из особенностей объекта (1.2.1) и задач анализа (1.2.2), в соответствии с таблицей 1.4.

\begin{gosttable}{}{Сопоставление групп средств по свойствам, существенным для навигации по большому дереву}

\begin{tabularx}{\linewidth}{>{\raggedright\arraybackslash}Xccccc}
\hline
Группа средств & Независимость от игры & Рост дерева во времени & Этапы итерации & Деревья свыше 10⁴ узлов & Автоматическая детализация \\
\hline
Учебные визуализаторы & − & ± & + & − & − \\
\hline
Универсальные библиотеки & + & − & − & ± & − \\
\hline
Средства просмотра итогового дерева & + & − & − & ± & − \\
\hline
Библиотеки с отображением & + & − & − & ± & − \\
\hline
Модуль предшествующей работы [1] & + & + & − & − & − \\
\hline
Просмотрщики профилей производительности [18] & + & − & − & + & + \\
\hline
\end{tabularx}

\end{gosttable}

Условные обозначения: «+» — свойство обеспечивается, «±» — обеспечивается частично, «−» — не обеспечивается.

Ни одно из рассмотренных средств не сочетает независимость от конкретной игры и наблюдение построения дерева с работоспособностью на больших деревьях и автоматическим управлением детализацией. Средства, обеспечивающие последние два свойства, — просмотрщики профилей производительности (1.4.5), — не предназначены для деревьев поиска: в них отсутствуют переиспользование дерева между ходами, различение объёма поиска и его результата, наблюдение роста и разбор итерации. Это сочетание и определяет незанятое место, которое должно закрыть разрабатываемое решение.

\subsection{Постановка задачи}

\subsubsection{Проблема}

Проведённый анализ позволяет сформулировать проблему, на решение которой направлена работа. Анализ работы метода поиска требует перехода между масштабами — от распределения усилий по всему дереву к статистике отдельного узла и обратно (1.2.2), однако число узлов дерева превосходит число одновременно различимых на экране элементов, а геометрическое уменьшение масштаба противоречия не разрешает (1.2.3). Существующие решения его также не снимают: разработанный ранее модуль [1] возлагает управление детализацией на исследователя и смещает узлы при сворачивании ветвей (1.3), а среди рассмотренных средств ни одно не сочетает независимости от конкретной игры и наблюдения за построением дерева с работоспособностью на больших деревьях и автоматическим управлением детализацией (1.5.4). Методы семантического масштабирования дают основу для устранения этого недостатка, однако их прямое применение затруднено особенностями дерева поиска (1.2.1) — неоднородностью распределения посещений, накоплением счётчиков вследствие переиспользования дерева между ходами игры и изменением дерева во времени.

Таким образом, проблема состоит в отсутствии способа визуализации дерева поиска, при котором уровень детализации согласуется с масштабом отображения без участия пользователя, представление сохраняет устойчивость как при изменении детализации, так и при росте дерева, а работоспособность не утрачивается при размере дерева в сотни тысяч узлов.

\subsubsection{Цель и исследовательские вопросы}

Целью работы является повышение наглядности и масштабируемости визуализации построения дерева поиска методом Монте-Карло в программной среде GGP Base Package за счёт применения методов семантического масштабирования.

Достижение цели предполагает ответ на четыре исследовательских вопроса:

\begin{enumerate}

  \item возможно ли согласовать уровень детализации с масштабом отображения так, чтобы число одновременно изображаемых узлов оставалось ограниченным при любом размере дерева и не требовало ручной настройки фильтров;

  \item возможно ли исключить смещение узлов при изменении детализации, то есть сохранить положение узла неизменным при догрузке и выгрузке соседних узлов;

  \item какое кодирование статистики узла обеспечивает одновременную различимость доминирующих и слаборазвитых ветвей дерева и сохраняет видимой границу его раскрытия;

  \item возможно ли сохранить пространственную устойчивость представления при наблюдении за ростом дерева, когда доли посещений изменяются с каждой итерацией поиска.

\end{enumerate}

Ответы на эти вопросы даются в последующих разделах и проверяются вычислительным экспериментом: первые два — измерением числа изображаемых узлов и смещения узлов при изменении масштаба в сопоставлении с раскладками, пересчитываемыми при каждой догрузке; третий — сопоставлением предлагаемого кодирования с его вариантами без нижнего порога доли и без разделения каналов; четвёртый — измерением смещения узлов между состояниями одного хода в сопоставлении с раскладкой по текущим долям посещений.

\subsubsection{Требования к решению и ограничения}

Базовым представлением выбрана ярусная диаграмма с горизонтальным семантическим масштабированием (1.4.7). Модель и реализующий её модуль должны удовлетворять следующим требованиям:

\begin{enumerate}

  \item отображать дерево поиска так, чтобы уровень детализации автоматически согласовывался с масштабом: на обзорном масштабе различимы доминирующие ветви, при приближении проявляется детализация выбранного поддерева;

  \item кодировать неоднородную статистику узлов так, чтобы одновременно были различимы как доминирующие, так и слаборазвитые ветви дерева;

  \item сохранять пространственную устойчивость представления — положение узлов не должно изменяться при изменении детализации, догрузке или выгрузке узлов;

  \item обеспечивать наблюдение за ростом дерева во времени без нарушения пространственной устойчивости;

  \item сохранять работоспособность и приблизительно постоянную плотность отображаемой информации при росте дерева от тысяч до сотен тысяч узлов и миллионах итераций поиска.

\end{enumerate}

Ограничения на решение таковы. Визуализация должна быть независимой от конкретной игры и опираться только на структуру дерева и статистику узлов, без игровых представлений. Доли посещений вычисляются относительно соседних узлов одного родителя, а не относительно корня, вследствие переиспользования дерева между ходами игры. Решение должно развивать модуль [1] и интегрироваться в среду GGP Base Package.

---

\subsection*{Список использованных источников}

\begin{enumerate}[label=\arabic*.]
  \item Ролдугин, О. Д. Создание в программной среде GGP Base Package модуля для визуализации работы метода поиска Монте-Карло по дереву игры : выпускная квалификационная работа бакалавра. — Волгоград : ВолгГТУ, 2025.
  \item Цыплаков, А. А. Пособие по теории игр / А. А. Цыплаков. — 2008. — 221 с.
  \item Genesereth, M., Thielscher, M. General Game Playing. — Morgan \& Claypool, 2014.
  \item Sironi, C. F. Monte-Carlo Tree Search for Artificial General Intelligence in Games : PhD thesis. — Maastricht University, 2019.
  \item Browne, C. B. et al. A Survey of Monte Carlo Tree Search Methods // IEEE Transactions on Computational Intelligence and AI in Games. — 2012. — Vol. 4, № 1. — P. 1–43.
  \item Chaslot, G. et al. Monte-Carlo Tree Search: A New Framework for Game AI // Proceedings of AIIDE. — 2008. — P. 216–217.
  \item Kocsis, L., Szepesvári, C. Bandit Based Monte-Carlo Planning // Machine Learning: ECML 2006. — Berlin, Heidelberg : Springer, 2006. — P. 282–293.
  \item Genesereth, M., Love, N., Pell, B. General Game Playing: Overview of the AAAI Competition // AI Magazine. — 2005. — Vol. 26, № 2. — P. 62–72.
  \item Shneiderman, B. The Eyes Have It: A Task by Data Type Taxonomy for Information Visualizations // Proceedings of the IEEE Symposium on Visual Languages. — 1996. — P. 336–343.
  \item Perlin, K., Fox, D. Pad: An Alternative Approach to the Computer Interface // SIGGRAPH '93. — 1993. — P. 57–64.
  \item Bederson, B. B., Hollan, J. D. Pad++: A Zooming Graphical Interface for Exploring Alternate Interface Physics // UIST '94. — 1994. — P. 17–26.
  \item Cockburn, A., Karlson, A., Bederson, B. B. A Review of Overview+Detail, Zooming, and Focus+Context Interfaces // ACM Computing Surveys. — 2009. — Vol. 41, № 1. — P. 1–31.
  \item Furnas, G. W. Generalized Fisheye Views // Proceedings of CHI '86. — 1986. — P. 16–23.
  \item Lamping, J., Rao, R., Pirolli, P. A Focus+Context Technique Based on Hyperbolic Geometry for Visualizing Large Hierarchies // Proceedings of CHI '95. — 1995. — P. 401–408.
  \item Plaisant, C., Grosjean, J., Bederson, B. B. SpaceTree: Supporting Exploration in Large Node Link Tree // Proceedings of IEEE InfoVis. — 2002. — P. 57–64.
  \item Heer, J., Card, S. K. DOITrees Revisited: Scalable, Space-Constrained Visualization of Hierarchical Data // Proceedings of AVI '04. — 2004. — P. 421–424.
  \item Kruskal, J. B., Landwehr, J. M. Icicle Plots: Better Displays for Hierarchical Clustering // The American Statistician. — 1983. — Vol. 37, № 2. — P. 162–168.
  \item Gregg, B. The Flame Graph // Communications of the ACM. — 2016. — Vol. 59, № 6. — P. 48–57.
  \item Shneiderman, B. Tree Visualization with Tree-Maps: A 2-D Space-Filling Approach // ACM Transactions on Graphics. — 1992. — Vol. 11, № 1. — P. 92–99.
  \item Stasko, J., Zhang, E. Focus+Context Display and Navigation Techniques for Enhancing Radial, Space-Filling Hierarchy Visualizations // Proceedings of IEEE InfoVis. — 2000. — P. 57–65.
  \item McMaster, R. B., Shea, K. S. Generalization in Digital Cartography. — Washington, DC : Association of American Geographers, 1992. — 134 p.
  \item Baier, H., Kaisers, M. Towards Explainable MCTS // AAAI Workshop on Explainable Agency in Artificial Intelligence. — 2021.
  \item Coulom, R. Treemaps for Search-Tree Visualization // Proceedings of the 7th Computer Olympiad Computer-Games Workshop. — Maastricht, 2002.
  \item Bruls, M., Huizing, K., van Wijk, J. J. Squarified Treemaps // Proceedings of Joint Eurographics and IEEE TCVG Symposium on Visualization. — Vienna : Springer, 2000. — P. 33–42.
  \item Elmqvist, N., Fekete, J.-D. Hierarchical Aggregation for Information Visualization: Overview, Techniques, and Design Guidelines // IEEE Transactions on Visualization and Computer Graphics. — 2010. — Vol. 16, № 3. — P. 439–454.
  \item Zinsmaier, M., Brandes, U., Deussen, O., Strobelt, H. Interactive Level-of-Detail Rendering of Large Graphs // IEEE Transactions on Visualization and Computer Graphics. — 2012. — Vol. 18, № 12. — P. 2486–2495.
  \item Ulmer, A., Angelini, M., Fekete, J.-D., Kohlhammer, J., May, T. A Survey on Progressive Visualization // IEEE Transactions on Visualization and Computer Graphics. — 2024. — Vol. 30, № 9. — P. 6447–6467.
  \item Archambault, D., Purchase, H. C., Pinaud, B. Animation, Small Multiples, and the Effect of Mental Map Preservation in Dynamic Graphs // IEEE Transactions on Visualization and Computer Graphics. — 2011. — Vol. 17, № 4. — P. 539–552.
  \item Литовкин, Д. В., Ролдугин, О. Д., Якимов, Г. А., Николаенко, С. Д., Смутин, Д. А. Динамическая визуализация построения MCTS-дерева в среде GGP Base Package // Моделирование, оптимизация и информационные технологии. — 2026. — Т. 14, № 7. — С. 1–11. — DOI: 10.26102/2310-6018/2026.58.7.001. — URL: https://moitvivt.ru/ru/journal/article?id=2340 (дата обращения: 14.11.2026).
  \item Garcia, V. mcts-viz: Visualization of MCTS Algorithm Applied to Tic-Tac-Toe [Электронный ресурс]. — URL: https://github.com/vgarciasc/mcts-viz (дата обращения: 20.03.2026).
  \item Varty, J. A Step-by-Step Look at Alpha Zero and Monte Carlo Tree Search [Электронный ресурс]. — URL: https://joshvarty.github.io/AlphaZero/ (дата обращения: 20.03.2026).
  \item AlphaZero Tic-Tac-Toe with Monte Carlo Tree Search [Электронный ресурс]. — URL: https://github.com/sagarnildass/AlphaZero-Tic-Tac-Toe-App (дата обращения: 20.03.2026).
  \item Gloqo: Game Logic Optimization Platform [Электронный ресурс]. — URL: https://www.gloqo.ai/visuals/mcts (дата обращения: 20.03.2026).
  \item D3: Data-Driven Documents [Электронный ресурс]. — URL: https://github.com/d3/d3 (дата обращения: 20.03.2026).
  \item Guo, B. [et al.] AlphaZero-Edu: Making AlphaZero Accessible to Everyone. — 2025. — 8 p. — DOI: 10.48550/arXiv.2504.14636.
  \item D3-Graphviz Integration Module [Электронный ресурс]. — URL: https://github.com/magjac/d3-graphviz (дата обращения: 20.03.2026).
  \item Zhao, Y., Hu, C., Liu, J. Playing With Monte-Carlo Tree Search [AI-eXplained] // IEEE Computational Intelligence Magazine. — 2024. — Vol. 19, № 1. — P. 85–86. — DOI: 10.1109/MCI.2023.3328150. — URL: https://aingames.cn/demo/mcts/index.html (дата обращения: 20.03.2026).
  \item General Game Playing: GitHub Repository [Электронный ресурс]. — URL: https://github.com/ggp-org/ggp-base (дата обращения: 20.03.2026).
  \item Json Grid Viewer Online [Электронный ресурс]. — URL: https://jsongrid.com/json-grid/ (дата обращения: 20.03.2026).
  \item MCTS.jl Julia Package [Электронный ресурс]. — URL: https://github.com/SISL/MCTS.jl (дата обращения: 20.03.2026).
  \item Вартанов, С. А. Введение в теорию игр : учебное пособие / С. А. Вартанов. — Москва : МАКС Пресс, 2018. — 140 с.
  \item Тихонова, О. В., Чихачева, О. А. Теория игр. Часть 1 : учебное пособие. — Рязань : Рязанский институт (филиал) Московского политехнического университета, 2020. — 48 с.
  \item Nash, J. F. Equilibrium Points in n-Person Games // Proceedings of the National Academy of Sciences. — 1950. — Vol. 36, № 1. — P. 48–49.
  \item Лепило, Н. Н. Теория игр : учебное пособие / Н. Н. Лепило. — Алчевск : ДонГТУ, 2020. — 132 с.
  \item Приложение: основы теории игр // Учебники экономического факультета МГУ [Электронный ресурс]. — URL: https://books.econ.msu.ru/Institutional-economics/app.GT/ (дата обращения: 29.05.2025).
\end{enumerate}
